\subsection{Time refinement of the same compact scalar detector}
\label{sec:source-scalar-time-refinement}

The golden tensor action admits a convergent local time discretization as
well as the spatial comparison above. This statement concerns the
mathematical trajectories of that action at each refinement. It does not
assert that their complete observer event histories have been executed.
Use the dimensionless coordinates, time and $\hbar=1$ of
\eqref{eq:common-scalar-action}. All norms below are in its mass inner
product, or equivalently in the mass-whitened coordinates.

\begin{proposition}[Full-spectrum displacement and detector estimate]
\label{prop:source-scalar-time-refinement}
Let $A$ be a finite positive self-adjoint operator, $A\ge I$, and use the
canonical kick--drift--kick evolution with step $\tau>0$, initial mean
position zero and initial mean momentum $v$. Let
$r=\tau^2\|A\|\le r_0<4$, $s_0=\sqrt{1-r_0/4}$ and
$d_0=s_0(1+s_0)$. At every integer $n\ge0$ with $n\tau\le T$,
\begin{equation}
 \|u_n-A^{-1/2}\sin(n\tau\sqrt A)v\|
 \le\tau^2\left[
 \frac{\|A^{1/2}v\|}{4d_0}+\frac{T\|Av\|}{12d_0}\right].
 \label{eq:source-scalar-time-displacement}
\end{equation}
Prepare the original action vacuum with coherent momentum displacement
$v$, and consider $E_g=(I+\sin\Phi(g))/2$. The probability difference
between this split evolution and continuous evolution of the same finite
action is at most
\begin{equation}
 \frac{\|g\|}{2}\tau^2\left[
 \frac{\|A^{1/2}v\|}{4d_0}+\frac{T\|Av\|}{12d_0}\right]
 +\frac{\tau^2\sqrt{G_0G_1}}{32(1-r_0/4)},
 \quad G_0=\|g\|^2,\quad G_1=\langle g,Ag\rangle.
 \label{eq:source-scalar-time-probability}
\end{equation}
Every mode is retained, including the split evolution of the original
vacuum covariance.
\end{proposition}
\begin{proof}
For an eigenfrequency $\omega>0$, put $z=\tau\omega/2$,
$s=\sqrt{1-z^2}$ and $\theta=2\arcsin z$.
The local canonical gates give the position recurrence
$u_{n+1}=2u_n-u_{n-1}-\tau^2Au_n$, with
$u_0=0$, $u_1=\tau v$. In this mode its solution is
$u_n=\sin(n\theta)v/(\omega s)$. Split its difference from the
continuous solution into amplitude and phase errors. Since
\[
 \frac{1/s-1}{\omega}=\frac{\tau^2\omega}{4s(1+s)},\qquad
 \theta-\tau\omega
 =2\int_0^z\frac{a^2\,da}{\sqrt{1-a^2}(1+\sqrt{1-a^2})}
 \le\frac{\tau^3\omega^3}{12d_0},
\]
the scalar error is at most
$\tau^2[\omega/(4d_0)+T\omega^2/(12d_0)]|v|$.
Functional calculus and the norm triangle inequality prove
\eqref{eq:source-scalar-time-displacement}.

In the original vacuum the mode variances are $1/(2\omega)$ and
$\omega/2$. The $n$-step canonical matrix has position row
$(\cos(n\theta),\sin(n\theta)/(\omega s))$. Its extra position variance is
\[
 \delta V_n(\omega)=
 \frac{\tau^2\omega\sin^2(n\theta)}{8(1-\tau^2\omega^2/4)}\ge0.
\]
Therefore the smeared variance increase is at most
$\tau^2\langle g,A^{1/2}g\rangle/[8(1-r_0/4)]$
and hence at most $\tau^2\sqrt{G_0G_1}/[8(1-r_0/4)]$.
For mean $m$ and variance $V$, the effect probability is
$[1+e^{-V/2}\sin m]/2$. Its change is bounded by
$|\delta m|/2+|\delta V|/4$; Cauchy--Schwarz bounds the mean change
by $\|g\|$ times \eqref{eq:source-scalar-time-displacement}.
This uses the original vacuum, not the invariant vacuum of a modified
discrete Hamiltonian.
\end{proof}

\paragraph{A same-action continuum limit.}
Let $h=\min_i(x_{i+1}-x_i)$ be the \emph{minimum} adjacent golden
coordinate gap, including the fixed Dirichlet endpoint. The one-dimensional
edge form satisfies
\[
 z^*Kz=\sum_i\frac{|z_{i+1}-z_i|^2}{h_i}
 \le\frac4h\sum_{i=1}^{q-1}|z_i|^2,
 \qquad z^*Wz\ge h\sum_{i=1}^{q-1}|z_i|^2.
\]
Consequently $\|A_q\|\le\Lambda_q:=1+12/h^2$.
For bounded mass norms of $v_q,g_q$, the displacement and probability
bounds vanish if $\tau_q=o(h)$. Indeed,
$\|A_q^{1/2}v_q\|\le\sqrt{\Lambda_q}\|v_q\|$,
$\|A_qv_q\|\le\Lambda_q\|v_q\|$ and
$\sqrt{G_0G_1}\le\sqrt{\Lambda_q}\|g_q\|^2$.
This requires no uniform derivative norm of a discontinuous compact
smearing. A concrete choice is the largest dyadic $\tau_q\le h^2$.
For $h\le1/4$, it gives $\tau_q^2\Lambda_q\le h^4+12h^2<1$.

Use exactly the sampled $f_c,g_c$ of
Theorem~\ref{thm:source-common-scalar-response}, the same original vacuum,
and $v_q=4f_c|_S$. The parent Gram bound $\|G-I\|\le d_q$ and positive
mass weights imply
\[
 \|g_c|_S\|^2\le1+d_q,\qquad
 \|v_q\|^2\le16(1+d_q).
\]
Multiplication by either slab indicator cannot increase these norms.
For $r_0=1$, $d_0\ge3/2$. On $0\le n\tau_q\le1$, the temporal
probability bound is therefore
\begin{equation}
 \mathcal E_{q,\mathrm{time}}
 \le\tau_q^2(1+d_q)
       \left[\frac38\sqrt{\Lambda_q}+\frac19\Lambda_q\right].
 \label{eq:source-scalar-time-compact-bound}
\end{equation}
Combined with Proposition~\ref{prop:common-scalar-continuum-limit}, this
proves convergence of the same compact displacement response and bounded
smeared field effect for the chosen space--time refinement. Bounded,
piecewise-continuous smearings whose products are Riemann integrable satisfy
the spatial hypothesis. No velocity-norm convergence, full-Fock state-norm
convergence across the growing lattice, or interacting field limit is
inferred from this estimate.

\begin{proposition}[A resolved fully time-discrete comparison]
\label{prop:source-scalar-time-resolved}
For the same $q=233$ source population, action, preparation and detector,
the exact minimum gap is $233-144\phi$ and the declared dyadic step is
$\tau=2^{-17}$. At every grid time $19/20\le n\tau\le1$, the
mathematical split evolution satisfies
\[
 |P^{\mathrm{split}}_{233,c}(n\tau)-P_c(n\tau)|<0.050063,
 \qquad |P^{\mathrm{split}}_{233,c}(n\tau)-\tfrac12|>0.137746.
\]
For every reference time $t\in[19/20,1]$, define the step readout by
\[
 k_0=\lceil(19/20)/\tau\rceil,\quad k_1=1/\tau,\quad
 k(t)=\min\{k_1,\max\{k_0,\lfloor t/\tau+1/2\rfloor\}\},
\]
\[
 P^{\mathrm{step}}_{233,c}(t)=P^{\mathrm{split}}_{233,c}(k(t)\tau).
\]
This is the nearest update inside the window, with upward ties.
Then $k(t)\tau\in[19/20,1]$, $|k(t)\tau-t|\le\tau$, and
\[
 |P^{\mathrm{step}}_{233,c}(t)-P_c(t)|<0.050078,\qquad
 |P^{\mathrm{step}}_{233,c}(t)-\tfrac12|>0.137746.
\]
This readout selects a defined discrete state; it introduces no state at
an intermediate time.
\end{proposition}
\begin{proof}
The separate exact certificate reconstructs the ordered golden coordinates
and minimum gap, verifies the dyadic step and stability inequality, and
replays the immutable spatial comparison. It bounds
\eqref{eq:source-scalar-time-compact-bound} by $0.000008071512$.
Adding this to the spatial error and subtracting it from the spatial
response gives the grid-time assertions. In the continuum the vacuum
variance is stationary and
$|P_c'(t)|\le\|g_c\|\|4f_c\|/2\le2$.
Rounding selects a nearest grid point before imposing the window, with
distance at most $\tau/2$. The upper endpoint $1$ is itself a grid point,
so upper clipping cannot occur for these $t$. If rounding falls below
$k_0$, clipping chooses $k_0\tau$, and
$(k_0-1)\tau<19/20\le t\le k_0\tau$ gives distance less than $\tau$.
Thus the selected point stays inside the window, also at its endpoints.
The extra comparison error is at most $2\tau$, and the discrete response
lower bound is unchanged.
\end{proof}

\begin{proposition}[Pre-arrival detector limit]
\label{prop:source-scalar-prearrival}
Use the same massive Dirichlet action on the cube, zero initial mean
position, original vacuum, and coherent momentum $af$, with fixed real
$a$. Let $f,g$ be bounded, compactly supported $L^2$ functions satisfying
the sampling hypotheses of Proposition~\ref{prop:common-scalar-continuum-limit}.
Suppose their support distance is $d>0$. For every $T<d$, the corresponding
local split trajectories with $\tau_q=o(h_q)$ satisfy
\[
 \sup_{0\le n\tau_q\le T}
 \left|P^{\mathrm{split}}_{q,g}(n\tau_q)-\tfrac12\right|\longrightarrow0.
\]
The same conclusion holds for the step readout
$P^{\mathrm{split}}_{q,g}(\lceil t/\tau_q\rceil\tau_q)$, uniformly for
$0\le t\le T$.
\end{proposition}
\begin{proof}
The continuum mean solves $u_{tt}-\Delta u+u=0$, with $u(0)=0$,
$u_t(0)=af$. Its finite propagation speed is one. Indeed the local energy
$e=(|u_t|^2+|\nabla u|^2+|u|^2)/2$ satisfies
$\partial_t e=\nabla\!\cdot(u_t\nabla u)$ and
$|u_t\nabla u\cdot n|\le e$ on each unit normal $n$.
Integrating over a shrinking backward light cone gives zero energy when
its initial base misses the preparation. The fixed Dirichlet boundary
has zero energy flux; smooth approximation in the energy norm extends the
argument to these $L^2$ momentum data. Thus $\langle g,u(t)\rangle=0$
for $t<d$. The stationary continuum vacuum covariance then gives the
sine-effect probability exactly $1/2$ there.

The parent spatial comparison is uniform on every fixed finite time
interval. Positive Riemann mass quadrature bounds the sampled norms,
so \eqref{eq:source-scalar-time-probability} makes the additional split
error vanish uniformly on that interval. This proves the grid assertion.
For the ceiling readout choose $T<T'<d$; eventually
$\lceil t/\tau_q\rceil\tau_q\le T+\tau_q<T'$ for all $t\le T$,
and apply the same uniform estimate on $[0,T']$.
\end{proof}

This is a limit of operational detector responses for the declared
preparation class. It does not identify microscopic read ancestry with
continuum causal order or establish a point-event count law. The
pre-arrival assertion is asymptotic; the finite error certificate above
concerns only its stated comparison window and preparation.

The finite local canonical gates define these trajectories and full-Fock
states for each fixed $q$ and finite number of steps. The certificate uses
source gaps, scalar norm bounds and the spatial comparison, not a stored
$232^3$-oscillator history. It does not import the separate $q=5$ preparation,
clock-error envelope or observer event receipt. The time refinement,
protected population, routing, action, boundary and quantization remain
declared. Native acceptance, a physical clock, outcome attestation and a
count-volume interpretation of substeps are separate requirements.
